% Source: deep_learning_cost_functions.tex
\begin{notechapter}
\chapter{Empirical Distribution, Maximum Likelihood and Cost Functions}\label{ch:dl-cost-functions}

\section{Introduction}
In deep learning, understanding the empirical distribution, maximum likelihood estimation, and the choice of cost functions is crucial for designing and training effective models. This document provides a comprehensive explanation of these concepts, their mathematical foundations, and their practical implications.

\section{Empirical Distribution}
The \textbf{empirical distribution} \( \hat{p}_{\text{data}} \) is a probability distribution derived directly from the observed training data. It represents the frequency of occurrences of different data points in the dataset.

\subsection{Definition}
Given a dataset \( \mathcal{D} = \{x^{(1)}, x^{(2)}, \dots, x^{(N)}\} \), the empirical distribution is defined as:
\[
\hat{p}_{\text{data}}(x) = \frac{1}{N} \sum_{i=1}^N \delta(x - x^{(i)})
\]
where \( \delta(x - x^{(i)}) \) is the Dirac delta ``function'' centered at \( x^{(i)} \): a unit point mass, zero everywhere except at \( x = x^{(i)} \) and integrating to \( 1 \) (it is not a function with value \( 1 \) at \( x^{(i)} \), but a distribution, defined by \( \int g(x)\,\delta(x - x^{(i)})\,dx = g(x^{(i)}) \)). For discrete \( x \), the empirical distribution simply assigns probability \( 1/N \) to each data point (with multiplicity), i.e.\ \( \delta \) is replaced by the indicator \( \mathbb{I}(x = x^{(i)}) \).

\subsection{Interpretation}
The empirical distribution treats each data point in the training set as equally likely. It essentially "counts" how often each value appears in the dataset and normalizes it by the total number of data points. For example, if a specific value \( x \) appears \( k \) times in the dataset, its probability under the empirical distribution is \( \frac{k}{N} \).

\subsection{Importance in Deep Learning}
\begin{itemize}
    \item \textbf{Learning Objective}: In supervised learning, the goal is often to minimize the difference between the model's predicted distribution \( p_{\text{model}}(x) \) and the empirical distribution \( \hat{p}_{\text{data}}(x) \). This is typically done using a loss function like cross-entropy.
    \item \textbf{Data Representation}: The empirical distribution serves as a proxy for the true underlying data distribution \( p_{\text{data}}(x) \), which is usually unknown.
    \item \textbf{Training Process}: During training, the model is optimized to match the empirical distribution as closely as possible, assuming that the training data is representative of the true distribution.
\end{itemize}

\section{Maximum Likelihood Estimation}
The goal of \textbf{maximum likelihood estimation (MLE)} is to find the model parameters \( \theta \) that make the observed data most probable under the model. In other words, we want to maximize the likelihood of the data given the model.

\subsection{Likelihood in Supervised Learning}
In supervised learning, the dataset \( \mathcal{D} \) consists of input-output pairs \( (x^{(i)}, y^{(i)}) \). The likelihood of the data under the model is the probability of the entire dataset given the parameters \( \theta \):
\[
p_{\text{model}}(\mathcal{D} | \theta) = \prod_{i=1}^N p_{\text{model}}((x^{(i)}, y^{(i)}) | \theta)
\]

\subsection{Factorizing the Joint Probability}
The joint probability \( p_{\text{model}}((x^{(i)}, y^{(i)}) | \theta) \) can be factorized using the chain rule of probability:
\[
p_{\text{model}}((x^{(i)}, y^{(i)}) | \theta) = p_{\text{model}}(x^{(i)} | \theta) \cdot p_{\text{model}}(y^{(i)} | x^{(i)}, \theta)
\]

\subsection{Conditional Likelihood}
In many supervised learning tasks, we are primarily interested in modeling the conditional distribution \( p_{\text{model}}(y^{(i)} | x^{(i)}, \theta) \), rather than the joint distribution. Thus, we often focus on the \textbf{conditional likelihood}:
\[
p_{\text{model}}(\mathcal{D} | \theta) = \prod_{i=1}^N p_{\text{model}}(y^{(i)} | x^{(i)}, \theta)
\]

\section{Negative Log-Likelihood as Cost Function}
The cost function in many modern neural networks is the \textbf{negative log-likelihood} because it aligns with the principle of maximum likelihood estimation.

\subsection{Log-Likelihood}
Working with products of probabilities can lead to numerical instability, so we take the logarithm of the likelihood, which turns the product into a sum:
\[
\log p_{\text{model}}(\mathcal{D} | \theta) = \sum_{i=1}^N \log p_{\text{model}}(y^{(i)} | x^{(i)}, \theta)
\]

\subsection{Negative Log-Likelihood}
In optimization, it's more common to minimize a cost function rather than maximize an objective. Therefore, we take the \textbf{negative log-likelihood} as the cost function:
\[
J(\theta) = -\sum_{i=1}^N \log p_{\text{model}}(y^{(i)} | x^{(i)}, \theta)
\]

\subsection{Expectation Form}
In practice, we often express the cost function as an expectation over the empirical distribution \( \hat{p}_{\text{data}} \):
\[
J(\theta) = -\mathbb{E}_{x,y \sim \hat{p}_{\text{data}}} \left[ \log p_{\text{model}}(y | x, \theta) \right]
\]

\subsection{Connection to Cross-Entropy}
The negative log-likelihood is equivalent to the \textbf{cross-entropy} between the empirical distribution \( \hat{p}_{\text{data}} \) and the model distribution \( p_{\text{model}} \). Cross-entropy measures the difference between two probability distributions:
\[
H(\hat{p}_{\text{data}}, p_{\text{model}}) = -\mathbb{E}_{x,y \sim \hat{p}_{\text{data}}} \left[ \log p_{\text{model}}(y | x, \theta) \right]
\]

\section{Connection Between Maximum Likelihood and KL Divergence}
The principle of maximum likelihood estimation is closely related to minimizing the \textbf{Kullback-Leibler (KL) divergence} between the empirical distribution \( \hat{p}_{\text{data}} \) and the model distribution \( p_{\text{model}} \).

\subsection{KL Divergence}
The KL divergence measures how one probability distribution diverges from a second, reference distribution:
\[
D_{\text{KL}}(P \| Q) = \mathbb{E}_{x \sim P} \left[ \log \frac{P(x)}{Q(x)} \right]
\]
In integral form, this becomes:
\[
D_{\text{KL}}(P \| Q) = \int P(x) \log \frac{P(x)}{Q(x)} \, dx
\]

\subsection{Forward KL Divergence}
The forward KL divergence \( D_{\text{KL}}(\hat{p}_{\text{data}} \| p_{\text{model}}) \) is:
\[
D_{\text{KL}}(\hat{p}_{\text{data}} \| p_{\text{model}}) = \mathbb{E}_{x,y \sim \hat{p}_{\text{data}}} \left[ \log \frac{\hat{p}_{\text{data}}(y | x)}{p_{\text{model}}(y | x, \theta)} \right]
\]
In integral form, this becomes:
\[
D_{\text{KL}}(\hat{p}_{\text{data}} \| p_{\text{model}}) = \int \hat{p}_{\text{data}}(y | x) \log \frac{\hat{p}_{\text{data}}(y | x)}{p_{\text{model}}(y | x, \theta)} \, dy \, dx
\]
This is the divergence minimized in maximum likelihood estimation.

\subsection{Reverse KL Divergence}
The reverse KL divergence \( D_{\text{KL}}(p_{\text{model}} \| \hat{p}_{\text{data}}) \) is:
\[
D_{\text{KL}}(p_{\text{model}} \| \hat{p}_{\text{data}}) = \mathbb{E}_{x,y \sim p_{\text{model}}} \left[ \log \frac{p_{\text{model}}(y | x, \theta)}{\hat{p}_{\text{data}}(y | x)} \right]
\]
In integral form, this becomes:
\[
D_{\text{KL}}(p_{\text{model}} \| \hat{p}_{\text{data}}) = \int p_{\text{model}}(y | x, \theta) \log \frac{p_{\text{model}}(y | x, \theta)}{\hat{p}_{\text{data}}(y | x)} \, dy \, dx
\]
This divergence is often used in variational inference and generative modeling.
Note however that the reverse KL divergence to the \emph{empirical} distribution is degenerate: as soon as \( p_{\text{model}} \) puts positive mass where \( \hat{p}_{\text{data}} = 0 \) (i.e.\ anywhere except exactly on the observed points), the ratio \( p_{\text{model}}/\hat{p}_{\text{data}} \) is infinite and so is the divergence. In practice the reverse KL is used against a target \emph{density} that can be evaluated, e.g.\ \( D_{\text{KL}}(q(z) \,\|\, p(z \mid x)) \) between an approximate and the true posterior in variational inference.\\

Figure \ref{dl-cost-functions:fig:fw_kl_vs_bw_kl} below is a visualization borrowed from
\href{https://arxiv.org/pdf/1804.00140}{this} article.

\begin{figure}[ht!]
    \centering
    \includegraphics[width=450px]{forward_KL_vs_backward_KL.png}
    \caption{Forward KL versus Backward KL}
    \label{dl-cost-functions:fig:fw_kl_vs_bw_kl}
\end{figure}

\subsection{Key Differences Between Forward and Reverse KL Divergence}
The choice of KL divergence has significant implications for the behavior of the model:

\begin{itemize}
    \item \textbf{Forward KL (\( D_{\text{KL}}(\hat{p}_{\text{data}} \| p_{\text{model}}) \))}:
        \begin{itemize}
            \item Penalizes the model when \( p_{\text{model}} \) assigns \textbf{low probability to data points that are likely under \( \hat{p}_{\text{data}} \)}.
            \item Encourages the model to cover all modes of the empirical distribution \( \hat{p}_{\text{data}} \), even if this means assigning some probability to regions where \( \hat{p}_{\text{data}} \) has no support.
            \item Used in maximum likelihood estimation and supervised learning.
        \end{itemize}
    \item \textbf{Reverse KL (\( D_{\text{KL}}(p_{\text{model}} \| \hat{p}_{\text{data}}) \))}:
        \begin{itemize}
            \item Penalizes the model when \( p_{\text{model}} \) assigns \textbf{high probability to data points that are unlikely under \( \hat{p}_{\text{data}} \)}.
            \item Encourages the model to avoid assigning probability to regions where \( \hat{p}_{\text{data}} \) has no support, even if this means ignoring some modes of \( \hat{p}_{\text{data}} \).
            \item Used in variational inference and generative modeling (e.g., Variational Autoencoders).
        \end{itemize}
\end{itemize}

\subsection{Practical Implications}
\begin{itemize}
    \item \textbf{Forward KL (\( D_{\text{KL}}(\hat{p}_{\text{data}} \| p_{\text{model}}) \))}:
        \begin{itemize}
            \item Suitable for supervised learning, where the goal is to match the model distribution to the empirical distribution.
            \item Ensures the model assigns high probability to all observed data points.
        \end{itemize}
    \item \textbf{Reverse KL (\( D_{\text{KL}}(p_{\text{model}} \| \hat{p}_{\text{data}}) \))}:
        \begin{itemize}
            \item Suitable for generative modeling, where the goal is to ensure the model does not generate unrealistic samples.
            \item Tends to produce "mode-seeking" behavior, where the model focuses on the dominant modes of the data distribution and ignores rare or noisy modes.
        \end{itemize}
\end{itemize}

\subsection{Mathematical Connection to Maximum Likelihood}
The forward KL divergence is directly related to maximum likelihood estimation. Minimizing \( D_{\text{KL}}(\hat{p}_{\text{data}} \| p_{\text{model}}) \) is equivalent to maximizing the log-likelihood of the data under the model:
\[
D_{\text{KL}}(\hat{p}_{\text{data}} \| p_{\text{model}}) = \text{constant} - \mathbb{E}_{x,y \sim \hat{p}_{\text{data}}} \left[ \log p_{\text{model}}(y | x, \theta) \right]
\]
In integral form, this becomes:
\[
D_{\text{KL}}(\hat{p}_{\text{data}} \| p_{\text{model}}) = \text{constant} - \int \hat{p}_{\text{data}}(y | x) \log p_{\text{model}}(y | x, \theta) \, dy \, dx
\]
Thus, minimizing the forward KL divergence is equivalent to minimizing the negative log-likelihood.

\section{Maximum Likelihood for Specific Distributions}
Here, we derive the maximum likelihood objective for three common output distributions: Gaussian, Bernoulli, and Multinoulli with softmax.

\subsection{Gaussian Distribution}
For a Gaussian distribution, the model predicts both the mean \( \mu(x, \theta) \) and the variance \( \sigma^2(x, \theta) \). The likelihood for a single example is:
\[
p_{\text{model}}(y | x, \theta) = \frac{1}{\sqrt{2\pi\sigma^2(x, \theta)}} \exp\left(-\frac{(y - \mu(x, \theta))^2}{2\sigma^2(x, \theta)}\right)
\]
The negative log-likelihood for a single example is:
\[
-\log p_{\text{model}}(y | x, \theta) = \frac{1}{2} \log(2\pi\sigma^2(x, \theta)) + \frac{(y - \mu(x, \theta))^2}{2\sigma^2(x, \theta)}
\]
For a dataset of \( N \) examples, the total negative log-likelihood (cost function) is:
\[
J(\theta) = \sum_{i=1}^N \left[ \frac{1}{2} \log(2\pi\sigma^2(x^{(i)}, \theta)) + \frac{(y^{(i)} - \mu(x^{(i)}, \theta))^2}{2\sigma^2(x^{(i)}, \theta)} \right]
\]

\subsubsection{Special Cases}
\begin{itemize}
    \item \textbf{Fixed Variance}: If the variance \( \sigma^2 \) is assumed to be constant, the cost function simplifies to minimizing the mean squared error (MSE):
    \[
    J(\theta) = \frac{1}{N} \sum_{i=1}^N (y^{(i)} - \mu(x^{(i)}, \theta))^2
    \]
    \item \textbf{Learned Variance}: If the variance is predicted by the model, the cost function includes both the mean and variance terms.
\end{itemize}

\subsubsection{Practical Considerations}
\begin{itemize}
    \item \textbf{Numerical Stability}: Predict the log variance \( \log \sigma^2(x, \theta) \) to avoid numerical issues.
    \item \textbf{Uncertainty Estimation}: Predicting the variance allows the model to estimate uncertainty in its predictions.
\end{itemize}

\subsection{Bernoulli Distribution}
The Bernoulli distribution is a special case of the Multinoulli distribution with \( K = 2 \). It models a binary random variable \( y \in \{0, 1\} \). The likelihood for a single example is:
\[
p_{\text{model}}(y | x, \theta) = p(x, \theta)^y (1 - p(x, \theta))^{1 - y}
\]
The negative log-likelihood is:
\[
-\log p_{\text{model}}(y | x, \theta) = -y \log p(x, \theta) - (1 - y) \log (1 - p(x, \theta))
\]
This is the binary cross-entropy loss. The probability \( p(x, \theta) \) is typically modeled using the sigmoid function:
\[
p(x, \theta) = \sigma(z(x, \theta)) = \frac{1}{1 + e^{-z(x, \theta)}}
\]
where \( z(x, \theta) \) is the output of the model before applying the sigmoid.

\subsection{Multinoulli Distribution with Softmax}
The Multinoulli distribution generalizes the Bernoulli distribution to \( K \) classes. It models a categorical random variable \( y \in \{1, 2, \dots, K\} \). The likelihood for a single example is:
\[
p_{\text{model}}(y | x, \theta) = \prod_{k=1}^K p_k(x, \theta)^{\mathbb{I}(y = k)}
\]
The negative log-likelihood is:
\[
-\log p_{\text{model}}(y | x, \theta) = -\sum_{k=1}^K \mathbb{I}(y = k) \log p_k(x, \theta)
\]
This is the categorical cross-entropy loss. The probabilities \( p_k(x, \theta) \) are typically modeled using the softmax function:
\[
p_k(x, \theta) = \frac{e^{z_k(x, \theta)}}{\sum_{j=1}^K e^{z_j(x, \theta)}}
\]
where \( z_k(x, \theta) \) is the output of the model for class \( k \) before applying the softmax.

\subsubsection{Softmax Reduces to Sigmoid for \( K = 2 \)}
When \( K = 2 \), the softmax function reduces to the sigmoid function. To see this, consider the softmax probabilities for two classes:
\[
p_1(x, \theta) = \frac{e^{z_1(x, \theta)}}{e^{z_1(x, \theta)} + e^{z_2(x, \theta)}}
\]
\[
p_2(x, \theta) = \frac{e^{z_2(x, \theta)}}{e^{z_1(x, \theta)} + e^{z_2(x, \theta)}}
\]
If we set \( z_2(x, \theta) = 0 \) (which is often done in practice for binary classification), then:
\[
p_1(x, \theta) = \frac{e^{z_1(x, \theta)}}{e^{z_1(x, \theta)} + 1} = \frac{1}{1 + e^{-z_1(x, \theta)}} = \sigma(z_1(x, \theta))
\]
\[
p_2(x, \theta) = \frac{1}{e^{z_1(x, \theta)} + 1} = 1 - \sigma(z_1(x, \theta))
\]
Thus, the softmax function reduces to the sigmoid function when \( K = 2 \), and the Multinoulli distribution reduces to the Bernoulli distribution.

\section{Conclusion}
The empirical distribution, maximum likelihood estimation, and the negative log-likelihood cost function are fundamental concepts in deep learning. By understanding these principles, we can design models that effectively learn from data and make accurate predictions. Additionally, the connection between maximum likelihood and KL divergence, as well as the derivations for specific distributions, provides deeper insights into the theoretical foundations of modern neural networks.
\end{notechapter}
